\begin{cabstract}
面向辅助交互与便携式使用场景，运动想象脑机接口的系统实现不仅取决于离线条件下的解码性能，还同时受到采集链路可实现性、任务与方法适配性以及在线交互可用性的共同约束。现有证据仍主要建立在离线、标准协议和相对高密度导联条件之上，对于低导联采集、边缘部署和在线反馈约束下如何形成可用闭环，仍缺少连续的系统性研究。围绕上述问题，本文以双导便携式脑机接口原型为对象，沿“可用链路构建、可部署解码、可在线验证”的主线展开研究，形成了从硬件系统、解码模型到在线验证的系统化实现路线。

针对便携式使用条件下采集、传输与基础交互链路缺乏系统级实现的问题，本文设计并实现了以 KS1092 采集前端和 nRF52832 主控为核心的双导脑电采集与基础交互系统，构建了脑电采样、参数配置、BLE 无线传输、事件标记和上位机基础交互链路。标签链路测试与 Alpha 节律验证表明，该系统能够实现稳定采集与事件同步，并在现有供电条件下连续运行约 6\,h，为后续算法验证与在线实验提供了可复用的硬件系统载体。针对低导联采集与端侧资源受限条件下解码有效性、结构紧凑性和部署友好性难以兼顾的问题，本文提出了基于局部窗口自注意力的运动想象脑电解码模型 EdgeMIFormer，并通过标准输入、导联受限和端侧部署等实验对其适用性进行验证。在标准 22 导四分类实验中，EdgeMIFormer 取得 83.76\% 的平均最佳准确率；经结构化剪枝与 8-bit 量化后，模型大小压缩至 655.36\,KB，并在 RK NPU 上实现约 31\,ms 的单次推理。

针对系统进一步走向在线交互时模型快速适配、系统协同与反馈闭环验证不足的问题，本文构建了“预训练初始化、轮次级模型更新与跨轮次预测反馈”相结合的在线验证机制。伪在线冷启动实验表明，有限参数重拟合能够改善早期累计表现；双导在线闭环实验进一步打通了采集、记录、训练、推理与反馈链路，部分完整会话的在线准确率稳定在 75\% 以上。上述研究初步贯通了低通道运动想象脑机接口从“能采集”到“能解码”再到“能在线验证”的技术路线，可为低成本、便携式、端侧脑机接口在系统构建、模型设计、部署实现和在线验证方面提供参考。

\end{cabstract}

\ckeywords{运动想象脑电；便携式脑机接口；自注意力解码；端侧部署；在线验证}

\begin{eabstract}
For application-oriented and portable usage scenarios, the implementation of motor imagery brain-computer interface systems is constrained not only by offline decoding performance, but also by the feasibility of the acquisition chain, the adaptation between task design and decoding methodology, and the usability of online interaction. Existing evidence is still mainly established under offline settings, standard protocols, and relatively high-density montages. How to form a usable closed loop under low-channel acquisition, edge deployment, and online feedback constraints remains insufficiently studied. To address these issues, this thesis investigates a dual-channel portable BCI prototype along a system pathway of usable signal-chain construction, deployable decoding, and online validation, forming a systematic implementation route from hardware system design to decoding and online validation.

To address the lack of a system-level solution for acquisition, transmission, and basic interaction under portable usage conditions, a dual-channel EEG acquisition system centered on the KS1092 front-end and the nRF52832 controller was designed and implemented, establishing signal acquisition, parameter configuration, BLE transmission, event tagging, and host-side interaction. Marker-chain tests and alpha-rhythm validation show that the system supports stable acquisition and event synchronization, with an operating time of about 6\,h under the current power setup, thereby providing a real hardware carrier for subsequent algorithm validation and online experiments. To address the difficulty of simultaneously ensuring decoding effectiveness, compactness, and deployment friendliness under low-channel acquisition and constrained edge-side resources, EdgeMIFormer, a motor imagery EEG decoding model based on local-window self-attention, was proposed and verified through standard-input evaluation, channel-constrained experiments, and edge-side deployment. In standard 22-channel four-class experiments, EdgeMIFormer achieved an average best accuracy of 83.76\%. After structured pruning and 8-bit quantization, the model size was reduced to 655.36\,KB and the single-run inference latency was about 31\,ms on an RK NPU.

To address the lack of evidence for rapid adaptation, system coordination, and feedback closure when the system moves toward online interaction, an online validation mechanism combining pretrained initialization, round-wise model updating, and cross-round prediction feedback was constructed. Pseudo-online cold-start experiments indicate that lightweight parameter refitting improves early cumulative performance, while real dual-channel online runs establish a complete loop of acquisition, logging, training, inference, and feedback. In some complete sessions, the online accuracy remained stably above 75\%. These results preliminarily connect the pathway from acquisition to decoding and online validation in low-channel motor imagery BCIs, and offer practical evidence for system construction, model design, edge deployment, and online validation in low-cost portable edge-oriented BCI systems.

\end{eabstract}

\ekeywords{Motor imagery EEG; Portable brain-computer interface; Self-attention decoding; Edge deployment; Online validation}
